% Options for packages loaded elsewhere
\PassOptionsToPackage{unicode}{hyperref}
\PassOptionsToPackage{hyphens}{url}
\documentclass[
]{article}
\usepackage{xcolor}
\usepackage{amsmath,amssymb}
\setcounter{secnumdepth}{-\maxdimen} % remove section numbering
\usepackage{iftex}
\ifPDFTeX
  \usepackage[T1]{fontenc}
  \usepackage[utf8]{inputenc}
  \usepackage{textcomp} % provide euro and other symbols
\else % if luatex or xetex
  \usepackage{unicode-math} % this also loads fontspec
  \defaultfontfeatures{Scale=MatchLowercase}
  \defaultfontfeatures[\rmfamily]{Ligatures=TeX,Scale=1}
\fi
\usepackage{lmodern}
\ifPDFTeX\else
  % xetex/luatex font selection
\fi
% Use upquote if available, for straight quotes in verbatim environments
\IfFileExists{upquote.sty}{\usepackage{upquote}}{}
\IfFileExists{microtype.sty}{% use microtype if available
  \usepackage[]{microtype}
  \UseMicrotypeSet[protrusion]{basicmath} % disable protrusion for tt fonts
}{}
\makeatletter
\@ifundefined{KOMAClassName}{% if non-KOMA class
  \IfFileExists{parskip.sty}{%
    \usepackage{parskip}
  }{% else
    \setlength{\parindent}{0pt}
    \setlength{\parskip}{6pt plus 2pt minus 1pt}}
}{% if KOMA class
  \KOMAoptions{parskip=half}}
\makeatother
\setlength{\emergencystretch}{3em} % prevent overfull lines
\providecommand{\tightlist}{%
  \setlength{\itemsep}{0pt}\setlength{\parskip}{0pt}}
\usepackage{bookmark}
\IfFileExists{xurl.sty}{\usepackage{xurl}}{} % add URL line breaks if available
\urlstyle{same}
% fallback for those not using the hyperref driver hyperxmp:
\makeatletter
\@ifundefined{xmpquote}{\newcommand{\xmpquote}[1]{#1}}{}
\makeatother
\hypersetup{
  hidelinks,
  pdfcreator={LaTeX via pandoc}}

\author{}
\date{}

\begin{document}

{
\setcounter{tocdepth}{2}
\tableofcontents
}
\section{WebCoder --- Ghidul
profesorului}\label{webcoder--ghidul-profesorului}

\begin{quote}
Rezumat într-un paragraf: WebCoder oferă profesorilor bucla de clasă în
stil pbinfo --- clase cu auto-înscriere a elevilor, teme cu ferestre de
publicare și expirare flexibilă, mod examen cu lockdown, seturi de grile
din banca Bac, grile de progres per elev, hărți de greșeli, tabele de
efort, fișe de evaluare și exporturi CSV/PDF --- peste un evaluator
bilingv real. Ghidul parcurge întregul flux de clasă.
\end{quote}

\textbf{Cui se adresează:} profesori, coordonatori de cercuri, antrenori
de olimpiadă județeană. Perechea în engleză:
\href{../guide-en/TEACHER-GUIDE.md}{\texttt{docs/guide-en/TEACHER-GUIDE.md}}.

\begin{center}\rule{0.5\linewidth}{0.5pt}\end{center}

\subsection{1. Obține rolul de
TEACHER}\label{1-obux21bine-rolul-de-teacher}

Cere unui administrator să-ți upgrade contul la \textbf{TEACHER}
(verificare pe email de școală / programul de parteneriate;
\texttt{docs/RUNBOOK.md} §Wave C). Rolul deblochează \texttt{/teach},
crearea de clase, authoring-ul de teme și suprafețele de analytics.

\subsection{2. Creează o clasă}\label{2-creeazux103-o-clasux103}

\begin{enumerate}
\def\labelenumi{\arabic{enumi}.}
\tightlist
\item
  \texttt{/teach} → \textbf{Clasă nouă} → numește-o (ex. „IX-A
  Informatică").
\item
  Primești un \textbf{cod de înscriere + parolă de înscriere}. Elevii se
  înscriu singuri din profil (modelul anti-abuz pbinfo --- tu nu creezi
  niciodată conturi de elevi):

  \begin{itemize}
  \tightlist
  \item
    distribuie codul, elevii îl introduc pe \texttt{/classroom};
  \item
    sau distribuie \textbf{codul QR} (înscriere prietenoasă pe mobil);
  \item
    există și import în masă pe listă de utilizatori (cei inexistenți
    sunt raportați).
  \end{itemize}
\item
  Lista membrilor leagă profilurile publice și vederile de submissions
  scope-uite pe clasă.
\end{enumerate}

\subsection{3. Temele}\label{3-temele}

Creează o temă pe o clasă:

\begin{itemize}
\tightlist
\item
  \textbf{probleme} alese individual, din capitole de curriculum sau din
  orice track într-un click („transformă acest track în temă");
\item
  \textbf{fereastra de disponibilitate} (\texttt{disponibilă\ de\ la} →
  \texttt{termen\ limită}, opțional \texttt{până\ la\ întârziere});
\item
  opțional \textbf{mod examen}: cât un elev e într-o temă de examen,
  restul site-ului (forum, alte probleme, indicii) e blocat pentru el,
  cu politica de conduită/AI publicată;
\item
  \textbf{seturi de grile} din banca Bacalaureat (filtra pe
  disciplină/clasa/an Bac, opțiune de amestecare).
\end{itemize}

Expirarea flexibilă (semanticile pbinfo): după deadline elevii pot
trimite în continuare --- acele submissions pur și simplu nu se numără
pentru temă. Evaluatorul nu pedepsește învățarea; upsolving-ul devine
repetiție.

\subsection{4. Urmărește clasa}\label{4-urmux103reux219te-clasa}

Tab-urile paginii clasei:

\begin{itemize}
\tightlist
\item
  \textbf{Grila} --- matrice elevi × probleme cu scorul maxim /
  verdictul / încercările, agregate peste submissions-urile din
  fereastră; sertar per elev cu istoricul complet în clasă.
\item
  \textbf{Insight} --- heatmap-ul de greșeli al clasei (ce probleme
  blochează pe cine), rate de nereușită per problemă.
\item
  \textbf{Tabela de efort} --- clasament pe încercări + zile active +
  rezolvări (tabla „grind").
\item
  \textbf{Fișa de evaluare} --- PDF per elev: scorurile per problemă din
  fiecare temă, totaluri, marcaje de întârziere. Printabilă pentru
  catalogul electronic.
\item
  \textbf{PDF de examen} --- temele se exportă ca broșură printabilă cu
  enunțuri + puncte.
\end{itemize}

\subsection{5. Exporturi și note}\label{5-exporturi-ux219i-note}

\begin{itemize}
\tightlist
\item
  catalog CSV (scorul maxim per problemă per elev + totaluri) pentru
  orice catalog extern;
\item
  fișe de evaluare PDF per clasă;
\item
  „descarcă toate soluțiile" per temă (cerința \#1 a profesorilor pe
  orice evaluator) pentru review offline de plagiat (plus badge-urile
  consultative de similaritate ale platformei).
\end{itemize}

\subsection{6. Arborele de curriculum}\label{6-arborele-de-curriculum}

Axa de curriculum din \texttt{/problems} mapează programa națională
(clasele 5--12 → capitole → subiecte) cu contoare de probleme per nod.
Construirea temelor = alegerea capitolelor: filtrezi catalogul pe
capitolul tău, selectezi, atribui. Arborele e dată versionată (o nouă
ediție de programă adaugă o versiune de taxonomie, nu o migrație).

\subsection{7. Organizează un concurs de
clasă}\label{7-organizeazux103-un-concurs-de-clasux103}

Concursurile cu fereastră se potrivesc orarului școlar: o fereastră de
câteva zile, fiecare elev primește o sesiune cronometrată (countdown
personal). Creează-l din \texttt{/contests} (TEACHER/CREATOR pot crea
concursuri private), atașează problemele clasei, distribuie înscrierea.
Cu sau fără rating --- fără rating e implicitul pentru runde de
antrenament. Pagina de replay arată timeline-ul oricărui elev minut cu
minut pentru analiza post-mortem.

\subsection{8. Creează-ți propriile
probleme}\label{8-creeazux103-ux21bi-propriile-probleme}

Conturile CREATOR/TEAVOR pot author-a:

\begin{itemize}
\tightlist
\item
  \textbf{wizard-ul vizual} în 7 pași (baze → limite → exemple → teste →
  soluție de referință → traducere → review), cu linter-ul de enunț și
  scorul de lizibilitate;
\item
  teste ascunse prin DSL-ul de generatoare sau prin tabla de bounties
  (contribuția comunității se recompensează);
\item
  import de pachete Polygon/Kattis (poartă de ledger-ul de permisiuni).
\end{itemize}

Orice problemă trece prin DRAFT → PENDING → APPROVED cu feedback de la
verificator; authoring-ul ÎNSEAMNĂ dezvoltare profesională (pagina
programului de contributori documentează specificațiile și șabloanele de
feedback).

\subsection{9. Certificări și
parteneriate}\label{9-certificux103ri-ux219i-parteneriate}

\begin{itemize}
\tightlist
\item
  Certificarea profesorilor (pagina programului explică cerințele și
  examenul);
\item
  parteneriate cu școli (șablon DPA în \texttt{docs/DPA-TEMPLATE.md});
\item
  prelucrarea datelor: datele de clasă urmează secțiunea GDPR din
  paginile legale; consimțământul minorilor e mediat de școală (conturi
  create de profesor sub baza legală a școlii) sau parental.
\end{itemize}

\subsection{10. Întrebări
frecvente}\label{10-uxeentrebux103ri-frecvente}

\textbf{Un elev a uitat parola} --- resetare self-service pe email; fără
email, contul mediat de școală e modelul robust. \textbf{Pot vedea codul
elevilor?} Da --- vederi de submissions scope-uite pe clasă, per elev
(grila le leagă). \textbf{Ratingurile de concurs afectează notele?} Nu
--- ratingurile și scorurile temelor sunt sisteme separate; exporti ce
ai nevoie.

\end{document}
